\chapter{2016年上海卷（秋文）}

\section{填空题}

\begin{problem}
设 \(x\in\R\)，则不等式 \(|x-3|<1\) 的解集为\fillinblank{}. 

\end{problem}

\begin{answer}
\((2,4)\)
\end{answer}

\begin{solution}
由 \(|x-3|<1\)，得 \(-1<x-3<1\)，所以 \(2<x<4\).
\end{solution}
\begin{problem}
设 \(z=\dfrac{3+2\i}{\i}\)，其中 \(\i\) 为虚数单位，则 \(z\) 的虚部等于\fillinblank{}. 

\end{problem}

\begin{answer}
\(-3\)
\end{answer}

\begin{solution}
\(z=\frac{3+2\i}{\i}=2-3\i\)，故虚部为 \(-3\).
\end{solution}
\begin{problem}
已知平行直线 \(l_1:2x+y-1=0\)，\(l_2:2x+y+1=0\)，则 \(l_1\) 与 \(l_2\) 的距离是\fillinblank{}. 

\end{problem}

\begin{answer}
\(\dfrac2{\sqrt5}\)
\end{answer}

\begin{solution}
距离为 \(\frac{|(-1)-1|}{\sqrt{2^2+1^2}}=\frac2{\sqrt5}\).
\end{solution}
\begin{problem}
某次体检，\(5\) 位同学的身高（单位：米）分别为 \(1.72,1.78,1.80,1.69,1.76\)，则这组数据的中位数是\fillinblank{}（米）. 

\end{problem}

\begin{answer}
\(1.76\)
\end{answer}

\begin{solution}
排序后为 \(1.69,1.72,1.76,1.78,1.80\)，中位数为 \(1.76\).
\end{solution}
\begin{problem}
若函数 \(f(x)=4\sin x+a\cos x\) 的最大值为 \(5\)，则常数 \(a=\fillinblank{}\). 

\end{problem}

\begin{answer}
\(\pm3\)
\end{answer}

\begin{solution}
函数 \(4\sin x+a\cos x\) 的最大值为 \(\sqrt{16+a^2}\). 由 \(\sqrt{16+a^2}=5\)，得 \(a=\pm3\).
\end{solution}
\begin{problem}
已知点 \((3,9)\) 在函数 \(f(x)=1+a^x\) 的图像上，则 \(f(x)\) 的反函数 \(f^{-1}(x)=\fillinblank{}\). 

\end{problem}

\begin{answer}
\(\log_2(x-1)\)
\end{answer}

\begin{solution}
由 \(9=1+a^3\)，得 \(a=2\). 故 \(f(x)=1+2^x\)，其反函数为 \(f^{-1}(x)=\log_2(x-1)\).
\end{solution}
\begin{problem}
若 \(x,y\) 满足 \(x\ge0\)，\(y\ge0\)，\(y\ge x+1\)，则 \(x-2y\) 的最大值为\fillinblank{}. 

\end{problem}

\begin{answer}
\(-2\)
\end{answer}

\begin{solution}
因为 \(y\ge x+1\)，所以 \(x-2y\le x-2(x+1)=-x-2\le-2\). 当 \(x=0,y=1\) 时取等号，故最大值为 \(-2\).
\end{solution}
\begin{problem}
方程 \(3\sin x=1+\cos2x\) 在区间 \([0,2\pi]\) 上的解为\fillinblank{}. 

\end{problem}

\begin{answer}
\(\dfrac{\pi}{6},\dfrac{5\pi}{6}\)
\end{answer}

\begin{solution}
由 \(1+\cos2x=2-2\sin^2x\)，得 \(2\sin^2x+3\sin x-2=0\)，故 \(\sin x=\frac12\). 因此 \(x=\frac{\pi}{6}\) 或 \(x=\frac{5\pi}{6}\).
\end{solution}
\begin{problem}
在 \(\left(\sqrt[3]{x}-\dfrac2x\right)^n\) 的二项展开式中，所有项的二项式系数之和为 \(256\)，则常数项等于\fillinblank{}. 

\end{problem}

\begin{answer}
\(112\)
\end{answer}

\begin{solution}
由 \(2^n=256\)，得 \(n=8\). 通项为 \(T_{k+1}=\mathrm C_8^k(-2)^kx^{\frac{8-4k}{3}}\). 令 \(\frac{8-4k}{3}=0\)，得 \(k=2\). 常数项为 \(\mathrm C_8^2(-2)^2=112\).
\end{solution}
\begin{problem}
已知 \(\triangle ABC\) 的三边长分别为 \(3,5,7\)，则该三角形的外接圆半径等于\fillinblank{}. 

\end{problem}

\begin{answer}
\(\dfrac{7\sqrt3}{3}\)
\end{answer}

\begin{solution}
由海伦公式得三角形面积为 \(\frac{15\sqrt3}{4}\). 因此外接圆半径 \(R=\frac{3\cdot5\cdot7}{4\cdot\frac{15\sqrt3}{4}}=\frac{7\sqrt3}{3}\).
\end{solution}
\begin{problem}
某食堂规定，每份午餐可以在四种水果中任选两种，则甲、乙两同学各自所选的两种水果相同的概率为\fillinblank{}. 

\end{problem}

\begin{answer}
\(\dfrac16\)
\end{answer}

\begin{solution}
每名同学可选的水果组合共有 \(\mathrm C_4^2=6\) 种. 甲的选择确定后，乙与甲相同的概率为 \(\frac16\).
\end{solution}
\begin{problem}
如图，已知点 \(O(0,0)\)，\(A(1,0)\)，\(B(0,-1)\)，\(P\) 是曲线 \(y=\sqrt{1-x^2}\) 上一个动点，则 \(\overrightarrow{OP}\cdot\overrightarrow{BA}\) 的取值范围是\fillinblank{}. 

\bitmapfigure[max width=3.3cm,max totalheight=3.3cm]{img/2016/shanghai_liberal/q12_fig1.png}

\end{problem}

\begin{answer}
\([-1,\sqrt2]\)
\end{answer}

\begin{solution}
设 \(P(x,y)\)，则 \(x^2+y^2=1\)，\(y\ge0\). 因为 \(\overrightarrow{OP}=(x,y)\)，\(\overrightarrow{BA}=(1,1)\)，所以 \(\overrightarrow{OP}\cdot\overrightarrow{BA}=x+y\). 在上半圆上，\(x+y\) 的最大值为 \(\sqrt2\)，最小值为 \(-1\)，故取值范围为 \([-1,\sqrt2]\).
\end{solution}
\begin{problem}
设 \(a>0\)，\(b>0\)，若关于 \(x,y\) 的方程组
\[
\begin{cases}
ax+y=1,\\
x+by=1
\end{cases}
\]
无解，则 \(a+b\) 的取值范围是\fillinblank{}. 

\end{problem}

\begin{answer}
\((2,+\infty)\)
\end{answer}

\begin{solution}
无解需 \(ab=1\)，且两方程不重合. 因此 \(b=\frac1a\)，且 \(a\ne1\). 所以 \(a+b=a+\frac1a>2\)，取值范围为 \((2,+\infty)\).
\end{solution}
\begin{problem}
无穷数列 \(\{a_n\}\) 由 \(k\) 个不同的数组成，\(S_n\) 为 \(\{a_n\}\) 的前 \(n\) 项和. 若对任意的 \(n\in\N^*\)，\(S_n\in\{2,3\}\)，则 \(k\) 的最大值为\fillinblank{}. 

\end{problem}

\begin{answer}
\(4\)
\end{answer}

\begin{solution}
因为 \(S_0=0\)，\(S_1\in\{2,3\}\)，故 \(a_1\) 只能为 \(2\) 或 \(3\). 当 \(n\ge2\) 时，\(a_n=S_n-S_{n-1}\in\{-1,0,1\}\). 因此不同项最多为 \(4\) 个. 例如令 \(S_n\) 依次取 \(3,2,3,3,2,\ldots\)，则 \(a_n\) 可出现 \(3,-1,1,0\)，所以最大值为 \(4\).
\end{solution}
\section{单选题}

\begin{problem}
设 \(a\in\R\)，则“\(a>1\)”是“\(a^2>1\)”的（\quad）

\choices
    {充分非必要条件}
    {必要非充分条件}
    {充要条件}
    {既非充分也非必要条件}

\end{problem}

\begin{answer}
A
\end{answer}

\begin{solution}
由 \(a>1\) 可推出 \(a^2>1\)，但 \(a^2>1\) 还可能有 \(a<-1\)，故选 \(A\).
\end{solution}
\begin{problem}
如图，在正方体 \(ABCD-A_1B_1C_1D_1\) 中，\(E,F\) 分别为 \(BC,BB_1\) 的中点，则下列直线中与直线 \(EF\) 相交的是（\quad）

\bitmapfigure[max width=3.5cm,max totalheight=3.5cm]{img/2016/shanghai_liberal/q16_fig1.png}

\choices
    {直线 \(AA_1\)}
    {直线 \(A_1B_1\)}
    {直线 \(A_1D_1\)}
    {直线 \(B_1C_1\)}

\end{problem}

\begin{answer}
D
\end{answer}

\begin{solution}
直线 \(EF\) 位于平面 \(BCC_1B_1\) 内，且其延长线与直线 \(B_1C_1\) 相交. 其他三条直线与 \(EF\) 异面或平行，故选 \(D\).
\end{solution}
\begin{problem}
设 \(a\in\R\)，\(b\in[0,2\pi)\). 若对任意实数 \(x\) 都有 \(\sin\!\left(3x-\dfrac{\pi}{3}\right)=\sin(ax+b)\)，则满足条件的有序实数对 \((a,b)\) 的对数为（\quad）

\choices
    {\(1\)}
    {\(2\)}
    {\(3\)}
    {\(4\)}

\end{problem}

\begin{answer}
B
\end{answer}

\begin{solution}
若两正弦函数恒等，则 \(ax+b=3x-\frac{\pi}{3}+2k\pi\)，或 \(ax+b=\pi-\left(3x-\frac{\pi}{3}\right)+2k\pi\). 结合 \(b\in[0,2\pi)\)，得 \((a,b)=(3,\frac{5\pi}{3})\) 或 \((-3,\frac{4\pi}{3})\)，共 \(2\) 对. 故选 \(B\).
\end{solution}
\begin{problem}
设 \(f(x),g(x),h(x)\) 是定义域为 \(\R\) 的三个函数. 对于命题：\(\circled{1}\) 若 \(f(x)+g(x)\)，\(f(x)+h(x)\)，\(g(x)+h(x)\) 均为增函数，则 \(f(x),g(x),h(x)\) 中至少有一个为增函数；\(\circled{2}\) 若 \(f(x)+g(x)\)，\(f(x)+h(x)\)，\(g(x)+h(x)\) 均是以 \(T\) 为周期的函数，则 \(f(x),g(x),h(x)\) 均是以 \(T\) 为周期的函数. 下列判断正确的是（\quad）

\choices
    {\(\circled{1}\) 和 \(\circled{2}\) 均为真命题}
    {\(\circled{1}\) 和 \(\circled{2}\) 均为假命题}
    {\(\circled{1}\) 为真命题，\(\circled{2}\) 为假命题}
    {\(\circled{1}\) 为假命题，\(\circled{2}\) 为真命题}

\end{problem}

\begin{answer}
D
\end{answer}

\begin{solution}
\(\circled{1}\) 为假命题，可取某段上三个函数增量为 \(-1,2,2\)，则两两和的增量均为正，但第一个函数并不递增. \(\circled{2}\) 为真命题，因为三组和函数均以 \(T\) 为周期，由 \(f+g\) 与 \(f+h\) 的周期性相减可得 \(g(x+T)-h(x+T)=g(x)-h(x)\)，再结合 \(g+h\) 的周期性可得 \(2g(x+T)=2g(x)\)，故 \(g\) 以 \(T\) 为周期. 继续由 \(f+g\)、\(g+h\) 的周期性分别推出 \(f,h\) 也以 \(T\) 为周期. 故选 \(D\).
\end{solution}
\section{解答题}

\begin{problem}

将边长为 \(1\) 的正方形 \(AA_1O_1O\)（及其内部）绕 \(OO_1\) 旋转一周形成圆柱，如图，\(\widehat{AC}\) 长为 \(\dfrac{5\pi}{6}\)，\(\widehat{A_1B_1}\) 长为 \(\dfrac{\pi}{3}\)，其中 \(B_1\) 与 \(C\) 在平面 \(AA_1O_1O\) 的同侧. 

\begin{enumerate}
\item 求圆柱的体积与侧面积；
\item 求异面直线 \(O_1B_1\) 与 \(OC\) 所成的角的大小. 
\end{enumerate}

\bitmapfigure[max width=4.3cm,max totalheight=3.6cm]{img/2016/shanghai_liberal/q19_fig1.png}

\end{problem}

\begin{solution}
（1）圆柱的底面半径为 \(1\)，高为 \(1\)，所以体积为 \(\pi\)，侧面积为 \(2\pi\).

（2）建立空间直角坐标系，使 \(O(0,0,0)\)，\(O_1(0,0,1)\)，\(A(1,0,0)\)，\(A_1(1,0,1)\). 因为 \(\widehat{AC}=\frac{5\pi}{6}\)，\(\widehat{A_1B_1}=\frac{\pi}{3}\)，且 \(B_1\) 与 \(C\) 在同侧，所以可取 \(C\left(-\frac{\sqrt3}{2},\frac12,0\right)\)，\(B_1\left(\frac12,\frac{\sqrt3}{2},1\right)\). 于是 \(\overrightarrow{O_1B_1}=(\frac12,\frac{\sqrt3}{2},0)\)，\(\overrightarrow{OC}=(-\frac{\sqrt3}{2},\frac12,0)\)，两向量点积为 \(0\). 因此异面直线 \(O_1B_1\) 与 \(OC\) 所成角为 \(\frac{\pi}{2}\).
\end{solution}
\begin{problem}

有一块正方形菜地 \(EFGH\)，\(EH\) 所在直线是一条小河，收获的蔬菜可送到 \(F\) 点或河边运走. 于是，菜地分为两个区域 \(S_1\) 和 \(S_2\)，其中 \(S_1\) 中的蔬菜运到河边较近，\(S_2\) 中的蔬菜运到 \(F\) 点较近，而菜地内 \(S_1\) 和 \(S_2\) 的分界线 \(C\) 上的点到河边与到 \(F\) 点的距离相等. 现建立平面直角坐标系，其中原点 \(O\) 为 \(EF\) 的中点，点 \(F\) 的坐标为 \((1,0)\)，如图. 

\begin{enumerate}
\item 求菜地内的分界线 \(C\) 的方程；
\item 菜农从蔬菜运量估计出 \(S_1\) 面积是 \(S_2\) 面积的两倍，由此得到 \(S_1\) 面积的“经验值”为 \(\dfrac83\). 设 \(M\) 是 \(C\) 上纵坐标为 \(1\) 的点，请计算以 \(EH\) 为一边、另有一边过点 \(M\) 的矩形的面积，及五边形 \(EOMGH\) 的面积，并判断哪一个更接近于 \(S_1\) 面积的“经验值”. 
\end{enumerate}

\bitmapfigure[max width=4.3cm,max totalheight=4.1cm]{img/2016/shanghai_liberal/q20_fig1.png}

\end{problem}

\begin{solution}
（1）设分界线 \(C\) 上一点为 \((x,y)\). 点到河边 \(EH:x=-1\) 的距离为 \(x+1\)，点到 \(F(1,0)\) 的距离为 \(\sqrt{(x-1)^2+y^2}\)，所以 \(x+1=\sqrt{(x-1)^2+y^2}\). 化简得 \(y^2=4x\). 因此菜地内分界线 \(C\) 的方程为 \(y^2=4x\)（\(0\le y\le2\)）.

（2）由 \(M\) 的纵坐标为 \(1\)，得 \(M\left(\frac14,1\right)\). 以 \(EH\) 为一边、另一边过点 \(M\) 的矩形面积为 \(2\cdot\frac54=\frac52\). 五边形 \(EOMGH\) 的顶点可取 \(E(-1,0)\)，\(O(0,0)\)，\(M(\frac14,1)\)，\(G(1,2)\)，\(H(-1,2)\)，面积为 \(\frac{11}{4}\). 因为 \(\left|\frac52-\frac83\right|=\frac16\)，\(\left|\frac{11}{4}-\frac83\right|=\frac1{12}\)，所以五边形 \(EOMGH\) 的面积更接近经验值.
\end{solution}
\begin{problem}

双曲线 \(x^2-\dfrac{y^2}{b^2}=1\)（\(b>0\)）的左、右焦点分别为 \(F_1,F_2\)，直线 \(l\) 过 \(F_2\) 且与双曲线交于 \(A,B\) 两点. 

\begin{enumerate}
\item 若 \(l\) 的倾斜角为 \(\dfrac{\pi}{2}\)，\(\triangle F_1AB\) 是等边三角形，求双曲线的渐近线方程；
\item 设 \(b=\sqrt3\)，若 \(l\) 的斜率存在，且 \(|AB|=4\)，求 \(l\) 的斜率. 
\end{enumerate}

\end{problem}

\begin{solution}
（1）设 \(c=\sqrt{1+b^2}\)，则 \(F_1(-c,0)\)，\(F_2(c,0)\). 因为 \(l\) 的倾斜角为 \(\frac{\pi}{2}\)，所以 \(l:x=c\). 代入双曲线得 \(y=\pm b^2\)，于是 \(A(c,b^2)\)，\(B(c,-b^2)\). 由 \(\triangle F_1AB\) 是等边三角形，得 \(\sqrt{(2c)^2+b^4}=2b^2\)，即 \(4(1+b^2)+b^4=4b^4\)，解得 \(b^2=2\). 故双曲线的渐近线方程为 \(y=\pm\sqrt2x\).

（2）当 \(b=\sqrt3\) 时，\(F_2(2,0)\). 设 \(l:y=k(x-2)\)，联立 \(x^2-\frac{y^2}{3}=1\)，得 \(\left(1-\frac{k^2}{3}\right)x^2+\frac{4k^2}{3}x-\frac{4k^2}{3}-1=0\). 两交点横坐标差的绝对值为 \(\frac{6\sqrt{1+k^2}}{|3-k^2|}\)，故弦长 \(|AB|=\frac{6(1+k^2)}{|3-k^2|}\). 由 \(|AB|=4\)，得 \(k^2=\frac35\). 所以 \(l\) 的斜率为 \(\pm\sqrt{\frac35}\).
\end{solution}
\begin{problem}

对于无穷数列 \(\{a_n\}\) 与 \(\{b_n\}\)，记 \(A=\{x\mid x=a_n,\ n\in\N^*\}\)，\(B=\{x\mid x=b_n,\ n\in\N^*\}\). 若同时满足条件：\(\circled{1}\) \(\{a_n\}\)，\(\{b_n\}\) 均单调递增；\(\circled{2}\) \(A\cap B=\varnothing\) 且 \(A\cup B=\N^*\)，则称 \(\{a_n\}\) 与 \(\{b_n\}\) 是无穷互补数列. 

\begin{enumerate}
\item 若 \(a_n=2n-1\)，\(b_n=4n-2\)，判断 \(\{a_n\}\) 与 \(\{b_n\}\) 是否为无穷互补数列，并说明理由；
\item 若 \(a_n=2^n\) 且 \(\{a_n\}\) 与 \(\{b_n\}\) 是无穷互补数列，求数列 \(\{b_n\}\) 的前 \(16\) 项的和；
\item 若 \(\{a_n\}\) 与 \(\{b_n\}\) 是无穷互补数列，\(\{a_n\}\) 为等差数列且 \(a_{16}=36\)，求 \(\{a_n\}\) 与 \(\{b_n\}\) 的通项公式. 
\end{enumerate}

\end{problem}

\begin{solution}
（1）不是无穷互补数列. 因为 \(a_n=2n-1\) 取全体正奇数，\(b_n=4n-2\) 取形如 \(4n-2\) 的正整数，故正整数 \(4,8,12,\ldots\) 既不属于 \(A\)，也不属于 \(B\)，即 \(A\cup B\ne\N^*\).

（2）因为 \(a_n=2^n\)，所以 \(A=\{2,4,8,16,\ldots\}\). 于是 \(\{b_n\}\) 的前 \(16\) 项为 \(1,3,5,6,7,9,10,11,12,13,14,15,17,18,19,20\)，其和为 \(\frac{20\cdot21}{2}-(2+4+8+16)=180\).

（3）设等差数列 \(a_n=a_1+(n-1)d\)，其中 \(a_1,d\in\N^*\). 由 \(a_{16}=36\)，得 \(a_1+15d=36\). 若 \(d=1\)，则 \(A\) 从 \(a_1\) 起包含所有充分大的正整数，补集有限，不可能与无穷数列 \(\{b_n\}\) 互补. 因此 \(d=2\)，\(a_1=6\). 故 \(a_n=2n+4\). 于是 \(B=\N^*\setminus A\)，所以
\[
b_n=
\begin{cases}
n,&1\le n\le5,\\
2n-3,&n\ge6.
\end{cases}
\]
\end{solution}
\begin{problem}

已知 \(a\in\R\)，函数 \(f(x)=\log_2\!\left(\dfrac1x+a\right)\). 

\begin{enumerate}
\item 当 \(a=1\) 时，解不等式 \(f(x)>1\)；
\item 若关于 \(x\) 的方程 \(f(x)+\log_2(x^2)=0\) 的解集中恰有一个元素，求 \(a\) 的值；
\item 设 \(a>0\)，若对任意 \(t\in\left[\dfrac12,1\right]\)，函数 \(f(x)\) 在区间 \([t,t+1]\) 上的最大值与最小值的差不超过 \(1\)，求 \(a\) 的取值范围. 
\end{enumerate}

\end{problem}

\begin{solution}
（1）当 \(a=1\) 时，\(f(x)>1\) 等价于 \(\frac1x+1>2\)，即 \(\frac1x>1\). 结合定义域可得解集为 \((0,1)\).

（2）由 \(f(x)+\log_2(x^2)=0\)，得 \(\log_2\left[\left(\frac1x+a\right)x^2\right]=0\)，即 \(ax^2+x-1=0\). 该方程的根自动满足原方程定义域. 当 \(a=0\) 时，方程为 \(x-1=0\)，有唯一解. 当 \(a\ne0\) 时，恰有一个实根需判别式 \(1+4a=0\)，得 \(a=-\frac14\). 因此 \(a=0\) 或 \(a=-\frac14\).

（3）当 \(a>0\) 时，\(f(x)\) 在 \((0,+\infty)\) 上单调递减. 因此对 \(t\in[\frac12,1]\)，最大值与最小值之差为 \(f(t)-f(t+1)\). 由 \(f(t)-f(t+1)\le1\)，得 \(\frac{1-t}{t(t+1)}\le a\). 函数 \(\frac{1-t}{t(t+1)}\) 在 \([\frac12,1]\) 上最大值为 \(\frac23\)，故 \(a\ge\frac23\).
\end{solution}
